\subsection{Recollections on FTL}

\begin{num}\label{num:multi-valued} \textit{Multi-valued series}.--
Formal ternary laws were introduced by Walter in an unpublished work,
 and then formally developped in \cite{DF3}. More recently, we further expanded the theory
 in \cite{CDFH}, introducing in particular a slightly more accurate framework that we now briefly recall,
 referring the reader to \cite[\textsection 1]{CDFH} for more details. 

 Let $R$ be a (commutative) ring. Given a pair of integers $(n,d) \in \ZZ^2$, $\ux=(x_1,...,x_d)$ a $d$-tuple of integers,
 an $n$-valued $d$-dimensional series, $(n,d)$-series for short,
 with coefficients in $R$ will be an element 
$$
F_t(\ux)=1+F_1(\ux)t+\hdots+F_n(\ux)t^n
$$
in $R[[\ux]][t]$, polynomial in $t$ of degree $n$.
 These $(n,d)$-series form an $R$-module $\MForm_{n,d}(R)$.
 We will use a \emph{$\ZZ$-filtration} on $\MForm_{n,d}(R)$ by assigning to the variables $x_i$ degree $+1$
 and to the variable $t$ degree $-1$. Therefore, an $(n,d)$-series has valuation greater or equal to $(-n)$
 and can have infinite degree.

One can define a generalized composition on these $(*,*)$-series
 (see \cite[Def. 2.1.13]{CDFH} for more details).
 Let $G_t(\uy)$ be an $(m,r)$-series. We say that $G_t(\uy)$ is \emph{composable}
 if $G_t(\underline 0)=1$. We then introduce formal variables $(G^{[1]},\hdots,G^{[m]})$,
 called the \emph{roots} of $G_t(\uy)$, satisfying the formal identity
\begin{equation}\label{eq:roots}
G_t(\uy)=\prod_{l=1}^m (1+G^{[l]}t).
\end{equation}
In other words, the coefficients of $G_t(\uy)$ are given by the elementary symmetric functions
 in the variables $G^{[l]}$.

Then we define the \emph{substitution of $G_t(\uy)$ in $F_t(\ux)$ at the variable $x_i$}
 as the $(nm,d+r-1)$-series in the variables $(x_1,\hdots,\cancel{x_i},\hdots,x_d,y_1,\hdots,y_r)$
 by the following equality:
$$
F_t\big(x_1,...,x_{i-1},G_t(\uy),x_{i+1},...,x_d\big)
 :=\prod_{l=1}^m F_t\big(x_1,...,x_{i-1},G^{[l]},x_{i+1},...,x_d\big).
$$
More precisely, the expression on the right is obviously symmetric in the variables $G^{[l]}$.
 Thus it can be expressed as a polynomial in the elementary symmetric polynomials $e_s$ in the
 $G^{[l]}$, and we can therefore substitute to the $e_s$ the actual coefficients of the
 polynomial $G_t(\uy)$, following the rule \eqref{eq:roots}.

The following algebraic notion is an analog of formal group laws,
 especially relevant in motivic homotopy theory (see also \cite[Def. 3.1.5]{CDFH}).
\end{num}
\begin{df}\label{df:FTL}
Let $R$ be a $\ZZe$-algebra.
A \emph{formal ternary law}, FTL for short,
 with coefficients in $R$ is a $(4,3)$-series with coefficients in $R$
$$
F_t(\ux)=1+F_1(x,y,z)t+F_2(x,y,z)t^2+F_3(x,y,z)t^3+F_4(x,y,z)t^4
$$
satisfying the following properties:
\begin{enumerate}
\item \textit{Neutral element}. The $(4,1)$-series $F_t(x,0,0)$ is split
with roots $x$ and $-\epsilon x$ each with multiplicity $2$,
 \emph{i.e.} $F_t(x,0,0)=(1+xt)^2(1-\epsilon xt)^2$.
\item \textit{Semi-neutral element}. The following relation holds:
$$
F_4(x,x,0)=0.
$$
\item \textit{Symmetry}. The element $F_t(x,y,z)$ of $R[[x,y,z]][t]$ is fixed
 under the action of the group $\mathfrak S(x,y,z)$ permuting the formal variables.
\item \textit{Associativity}. Given formal variables $(x,y,z,u,v)$,
 one has the following equality of $(16,5)$-series:
$$
F_t\big(F_t(x,y,z),u,v\big)=F_t\big(x,F_t(y,z,u),v\big)
$$
where we have used the above substitution operation as $F_t(x,y,z)$ is composable by property (1).
\item \textit{$\epsilon$-Linearity}. One has the following relation in $\MForm_{4,3}(R)$:
$$
F_t(-\epsilon x,y,z)=F_{-\epsilon t}(x,y,z).
$$
\end{enumerate}
We frequently display an FTL by its coefficients:
\begin{equation}\label{eq:generitc_FTL}
F_t(x,y,z)=1+\sum_{i,j,k\geq 0, 1 \leq l \leq 4} a_{ijk}^lx^iy^jz^kt^l.
\end{equation}
Then, according to the above paragraph, the \emph{degree} of $F_t(\ux)$ is the integer
$$
d=\max \{i+j+k-l \mid a_{ijk}^l \neq 0\}.
$$
\end{df}

\begin{rem} Most of the above axioms are easy to translate in terms of coefficients.
 Using the notation \eqref{eq:generitc_FTL}, one obtains the following translations:
\begin{enumerate}
\item \textit{Neutral element}:
 $a_{i00}^l=\begin{cases}
1 & i=l=4, \\
2(1-\epsilon) & i=l=1,3, \\
2(1-2\epsilon) & i=l=2, \\
0 & \text{otherwise.}
\end{cases}$
\item \textit{Semi-neutral element}: For all $n \geq 0$, $\sum_{i+j=n} a^4_{ij0}=0$.
\item \textit{Symmetry}: for all $l$, $a_{ijk}^l$ is invariant under permutations of the triple $(i,j,k)$.
\item[(5)] \textit{$\epsilon$-Linearity}. The relation $(1+\epsilon)a_{ijk}^l=0$ holds
 whenever $l$ and one of the $i,j,k$ do not have the same parity.
\end{enumerate}
The associativity axiom is computationaly very involved and we refer the reader to \cite[Appendix]{CDFH}
 for a computer-based calculation.
\end{rem}

\begin{ex}\label{ex:FTP_low_complexity}
 The following FTL are examples with bounded degree. As such, they are respective analogues
 of the additive and multiplicative formal group laws (FGL for short).
\begin{enumerate}
\item We have an FTL of \emph{degree $0$} with coefficients in $\ZZe$ whose non-zero coefficients are:
$$
\begin{array}{llll}
a^1_{100}=2(1-\epsilon) & & & \\
a^2_{200}=2(1-2\epsilon) & a^2_{110}=2(1-\epsilon) & & \\
a^3_{300}=2(1-\epsilon) & a^3_{210}=-2(1-\epsilon) & a^3_{111}=8(2-3\epsilon) & \\
a^4_{400}=1 & a^4_{310}=-2(1-\epsilon) & a^4_{220}=2(1-2\epsilon) & a^4_{211}=2(1-\epsilon).
\end{array}
$$
After inverting $2$, this is one of the two universal FTL of degree $0$, and even $\leq 0$: see \cite[Th. 3.1.12]{CDFH}.
 The other FTL of degree $0$ is given by the same coefficients except that $a^3_{111}=8(3-2\epsilon)$.
 Note that the universality statement tells us that these coefficients will appear in any FTL.
\item We have an FTL of \emph{degree $2$}, having parameters $\tau$ and $\gamma$,
 \emph{i.e.} with coefficients in the polynomial ring
 $\ZZe^{mul}:=\ZZe[\tau,\gamma^{\pm 1}]/\langle \tau^2-2(1-\epsilon)\gamma, (1+\epsilon)\tau\rangle$,
 whose non-zero coefficients are, in addition to the ones appearing in the preceding law:
$$
\begin{array}{llll}
a^1_{110}=\tau\gamma^{-1} & a^1_{111}=\gamma^{-1} & & \\
a^2_{210}=2\tau\gamma^{-1} & a^2_{111}=-3\tau\gamma^{-1} & a^2_{220}=\gamma^{-1} & \\
a^3_{310}=\tau\gamma^{-1} & a^2_{220}=-2\tau\gamma^{-1} & a^3_{211}=3\tau\gamma^{-1} & a^3_{220}=\gamma^{-1} \\
a^4_{311}=-\tau\gamma^{-1} & a^4_{221}=2\tau\gamma^{-1} & a^4_{222}=\gamma^{-1}. &
\end{array}
$$
\end{enumerate}
See \cite[Theorem 6.6]{FH21} for more details. Note that these two FTL have no non-trivial coefficients in degrees less than $0$.
\end{ex}

By analogy with the case of formal group laws as underlined above,
 we will adopt the following definition.
\begin{df}\label{df:addmulFTL}
The FTL with coefficients in $\ZZe$ (resp. $\ZZe^{mul}$) of point (1) of the above example
 will be called the \emph{additive} (resp. \emph{multiplicative}) FTL
 and denoted by $F_t^{add}$ (resp. $F_t^{mul}$).

We will use the same terminology for any FTL obtained by extension of scalars
 from the two previous ones.
\end{df}

\begin{num}
As for FGL, FTL can be organized in a category.
 First note that if $F_t(x,y,z)$ is an FTL with coefficients $a_{ijk}^l$ in a ring $R$,
 and $\varphi:R \rightarrow R'$ is a morphism of rings,
 one obviously obtains a new FTL with coefficients in $R'$ by applying $\varphi$ to all
 coefficients of $F_t(x,y,z)$. We will denote this new FTL by
$$
\varphi_*F_t(x,y,z):=1+\sum_{i,j,k\geq 0, 1 \leq l \leq 4} \varphi(a_{ijk}^l)x^iy^jz^kt^l.
$$
\end{num}
\begin{df}\label{df:morph_FTL}
Let $F_t(x,y,z)$ and $G_t(x,y,z)$ be FTL with coefficients respectively
 in $\ZZe$-algebras $R$ and $R'$. A morphism from $F_t(x,y,z)$ to $G_t(x,y,z)$ is pair $(\varphi,\Theta)$ where $\varphi:R \rightarrow R'$
 is a morphism of $\ZZe$-algebras,
 $\Theta \in R[[x]]$ is a composable power series such that
 the following equality of $(4,3)$-series holds:
$$
\Theta_t\big(\varphi_*F_t(x,y,z)\big)=G_t\big(\Theta_t(x),\Theta_t(y),\Theta_t(y)\big),
$$
where $\Theta_t(x)=1+\Theta(x).t$ as a $(1,1)$-series.
 Note that the two terms of this equality are well defined,
 using the substitution operation of \Cref{num:multi-valued},
 as both $\Theta_t(x)$ and $F_t(x,y,z)$ are composable. Explicitly, the above equality reads as
 \[
 \prod_{i=1}^4(1+\Theta(F^{[i]})t)=G_t\big(\Theta_t(x),\Theta_t(y),\Theta_t(y)\big).
 \]

The composition of these morphisms are defined by composition of rings for $\varphi$,
 and by composition of power series for $\Theta(x)$.
 Therefore an isomorphism of FTL is a pair $(\varphi,\Theta)$ such that $\varphi$ is an isomorphism
 of $\ZZe$-algebras, and $\Theta$ is an invertible power series.
 Such an isomorphism will be called \emph{strict} if $\theta(x)=x+\sum_{i\geq 2}a_ix^i$.

The corresponding category will be denoted by $\FTL$.
\end{df}

\begin{rem}
The category $\FTL$ is a $\ZZe$-linear category, cofibred over the category of $\ZZe$-algebras.
In particular, any morphism $(\varphi,\Theta)$ of FTL can be factorised as
 $(\varphi,\Theta)=(\Id,\Theta)\circ (\varphi,\Id)$.
We can then restrict to the case where FTL have the same coefficient ring.
\end{rem}

\begin{ex}
The \emph{$\epsilon$-linearity axiom} for an FTL $F_t(x,y,z)$ with coefficients in $R$ 
 can be translated by saying that the series $\Theta(x)=-\epsilon.x$, seen as a power series
 with coefficients in $R$,
 induces an automorphism $(\Id_R,\Theta(x))$ of $F_t(x,y,z)$.
\end{ex}

\begin{num}
Let us briefly recall the discussion of \cite[3.1.1]{CDFH}.
 We define the plus and minus parts of $\ZZe$ as the quotient
 rings $\ZZep=\ZZe/(1+\epsilon)$, $\ZZem=\ZZe/(1-\epsilon)$.
 There is an injective map $\ZZe \rightarrow \ZZep \times \ZZem$
 which is an isomorphism after inverting $2$.
 The same definition applies to an arbitrary $\ZZe$-algebra $R$,
 and $R \rightarrow R^+ \times R^-$ is injective if $R/\ZZe$ is flat
 and bijective if $2 \in R^\times$.

Consequently, given any FTL $(R,\F_t(x,y,z))$,
 one defines its plus and minus parts by base change 
 respectively to $R^+/R$ and $R^-/R$.
 Equivalently, the FTLs $\big(R^+,F_t^+(x,y,z)\big)$ and $\big(R^-,F_t^-(x,y,z)\big)$
 are respectively obtained by specializing the symbol $\epsilon$ to $(-1)$ and $(+1)$.
\end{num}

\begin{ex}
The minus part of the additive formal group law can be written explicitly:
$$
F_t^{add-}(x,y,z)=1-2x^2.t^2-8xyz.t^3+(x^4-2x^2y^2).t^4.
$$
\end{ex}


\subsection{$\Sp$-oriented theories}\label{sec:sp_orientaiton}

\begin{num}\label{num:sp-grassmanian}
In this section, we recall the links between $\Sp$-oriented ring spectra and symplectic cobordism, following Panin and Walter.

We first start with the construction of symplectic cobordism.
 Recall that Voevodsky introduced in \cite{VoeMilnor} the algebraic cobordism spectrum $\MGL_S$
 over a base scheme $S$, as the algebraic analogue of complex cobordism.
 As in the classical algebraic topological case,
 it is defined as the colimits of the Thom spaces of the universal rank $n$ vector bundle
 $\gamma_n$ on the grassmanian scheme $\Gr_n$ of $n$-vector subspaces of the infinite dimension affine
 space:
$$
\MGL_S=\colim_{n>0} \Omega^n \Th(\gamma_n).
$$
This is a (commutative) ring spectrum, in the strict sense
 (see \cite{PPR} for the construction in terms of symmetric spectra,
 and \cite[\textsection 16]{BH} for that in terms of monoidal $\infty$-categories).
 Recall lastly that thanks to a theorem of Morel and Voevodsky,
 $\Gr_n$ is weakly equivalent in the $\AA^1$-homotopy category to the classifying space $\BGL_n$
 of the general linear group $\GL_n$.

Panin and Walter introduced in \cite{PWcobord} the analogous construction
 by replacing the family $\GL_n$ by that of the symplectic group $\Sp_{2n}$,
 following again classical algebraic topology. The main difference is that
 one naturally produces a $(\PP^1)^{\wedge 2}$-spectrum.

First, the classifying space $\BSp_{2n}$ also admits a geometrical model,
 the hyperbolic (or ``quaternionic'') Grassmannian.
 To fix notations, we consider the standard symplectic form $h$ on $\AA^2_S$,
 and denote by $\omega_2$ the corresponding symplectic bundle.
 We also set $\omega_{2n}=\omega_{2n-2}\perp \omega_2$ for any $n\geq 2$.
 For any integer $0\leq r\leq n$, we can consider the \emph{hyperbolic Grassmannian} $\HGr_S(2r,\omega_{2n})$ defined to be the open subscheme of the Grassmannian $\Gr_S(2r,2n)$ of $2r$-subvector bundles of $\AA^{2n}$
 on which the restriction of $\omega_{2n}$ is non degenerate.
 It satisfies the following compatibilities:
\begin{enumerate}
\item We have a canonical embedding $\HGr_S(2r,\omega_{2n})\to \HGr_S(2r,\omega_{2n+2})$
 which classifies flags of vector bundles of the form $U_{2r}\subset p^*\omega_{2n}\subset  p^*\omega_{2n+2}$.
\item We have a canonical embedding $\HGr_S(2r,\omega_{2n})\to \HGr_S(2r+2,\omega_{2n+2})$ which classifies flags of vector bundles
 $U_{2r}\perp p^*\omega_2\subset p^*\omega_{2n}\perp p^*\omega_{2}$.
\end{enumerate}

Then we get an $\AA^1$-weak equivalence of pointed spaces:
$$
\BSp_{2r}=\colim_{n>0} \HGr_S(2r,\omega_{2n}).
$$
We also consider the \emph{hyperbolic projective space} $\HP^n_S=\HGr_S(2,\omega_{2n+2})$,
 of dimension $2n$ over $S$, and $\HP^\infty_S=\colim_{n>0} \HGr_S(2,\omega_{2n+2})$,
 so that $\HP^\infty_S=\BSp_2$. Note that $\HP^1_S \simeq (\PP^1)^{\wedge 2}$, the two-fold smash-power
 of the Tate sphere.

We then define $\MSp_{2n}=\Th(\cU_{2n})$, as a space over $S$,
 $\cU_{2r}$ being the tautological symplectic bundle on $\BSp_{2n}$.
 The second of the above properties allows to define $(\PP^1)^{\wedge 2}$-suspension maps as follows.
 The embedding $f:\HGr_S(2n,\omega_{2n})\to \HGr_S(2n+2,\omega_{2n+2})$ yields a bundle of rank $2n+2$ on the first factor,
 namely $f^*\cU_{2n+2}$ which splits as $f^*\cU_{2n+2}\simeq \cU_{2n}\oplus \omega_2$.
 We obtain a morphism of Thom spaces
\[
\Th(\cU_{2n})\wedge th(\omega_2)\simeq \Th(\cU_{2n}\oplus \omega_2)\simeq \Th(f^*\cU_{2n+2})\to \Th(\cU_{2n+2})
\]
which is compatible with the inclusions of item 1. This yields the desired suspension maps:
\[
\sigma_{2n}:(\PP^1)^{\wedge 2} \wedge \MSp_{2n} \simeq \Th(\omega_2)\wedge \MSp_{2n}\to \MSp_{2n+2},
\]
and therefore a $(\PP^1)^{\wedge 2}$-spectrum, which in $\infty$-categorical terms can be defined as:
\begin{equation}\label{eq:MSp_as_colim}
\MSp_S=\colim_{r>0} \Omega^{2n} \MSp_{2n}.
\end{equation}
To get the ring structure, one considers the rank $2(n+m)$ symplectic bundle
 $\cU_{2n} \boxplus \cU_{2m}$ over $\BSp_{2n} \times_S \BSp_{2m}$,
 obtained by exterior sum. It is classified by a map
 $$f:\BSp_{2n} \times_S \BSp_{2m} \rightarrow \BSp_{2(n+m)}$$ such that
 $f^{-1}(\cU_{2(n+m)})=\cU_{2n} \boxplus \cU_{2m}$. Taking Thom spaces, one deduces:
$$
\Th(\cU_{2n}) \otimes \Th(\cU_{2m})
 \simeq \Th(\cU_{2n} \boxplus \cU_{2m}) \xrightarrow{f'_*} \Th(\cU_{2(n+m)})
$$
as required. This gives a (commutative) ring structure on the spectrum $\MSp_S$,
 which can be lifted to a strict ring structure at the level of symmetric $(\PP^1)^{\wedge 2}$-spectra
 according to \cite{PWcobord}.\footnote{For a more conceptual construction
 of the corresponding $E_\infty$-ring structure, we refer the reader to \cite[\textsection 16]{BH}.
 However, in this paper, we will only need the (commutative) monoid structure in $\SH(S)$
 on $\MSp_S$.}
\end{num}
\begin{df}
The ring spectrum $\MSp_S$ is called the symplectic cobordism spectrum over $S$.
\end{df}

\begin{num}\label{num:symplectic_Thom_Grass_Euler}
 \textit{Thom classes for symplectic cobordism}.
As in topology, the previous ring spectrum is built to possess ``symplectic Thom classes''.
 First note that the identity of $\MSp_{2r}$ corresponds by construction of $\MSp_S$
 to a map:
\[
\tau_r:\Sigma^{\infty}\MSp_{2r}\to \MSp(2r)[4r],
\]
or in other words, an element $\tau_r \in \tilde \MSp^{4r,2r}(\MSp_{2r})$
 in the reduced symplectic cobordism.
 Let $X$ be a smooth $S$-scheme which is affine.
 Then a symplectic bundle $\cE$ of rank $2r$ is classified by a map:
\[
f:X\to \BSp_{2r}
\] 
such that $\cE=f^{-1}(\cU_{2r})$. Taking Thom spaces and composing with $\tau_r$, we get:
\[
Th(\cE)=Th(f^{-1}\cU_{2r}) \xrightarrow{f_*} Th(\cU_{2r})=\MSp_{2r} \xrightarrow{\tau_r} \MSp(2r)[4r]
\]
This is the Thom class $\thom^\MSp(\cE) \in \tilde \MSp^{4r,2r}(\Th(\cE))$
 associated with $\cE$
 with coefficients in $\MSp$. Equivalently, $\thom(\cE)=f^*(\tau_r)$.

Using Jouanolou's trick, this construction can be extended to any smooth $S$-scheme.
 One can check the basic properties of Thom classes, as stated in \cite[Def. 2.1.3, $G=\Sp_*$]{DF3}:
 compatibility with isomorphisms and pullbacks, products, normalisation.
 In other words, $\MSp_S$ is symplectically oriented in the sense of \emph{loc. cit.}
 (see \cite[\textsection 10, Th. 1.1]{PWcobord} for the proof, analogous to the
 topological case).
%An $\MSp$-algebra over $S$ is a morphism of ring spectra $\varphi:\MSp_S \rightarrow \E$
% over $S$.
% We immediately associates to $\varphi$ Thom classes $\thom^\E$ for $\E$ by the formula:
%$$
%\thom^\E(\cU)=\varphi_*(\thom^\MSp(\cU)).
%$$
%One immediately check that this provides a structure of an $\Sp$-orientation on $\E$.
Let us introduce some terminology for the statement of the next result.
 Let $\E$ be a ring spectrum over $S$.
 When $\E$ admits Thom classes for symplectic bundles satisfying the properties
 mentioned above, we will say that $\E$ admits symplectic Thom classes --- this corresponds
 to being $\Sp$-oriented in the sense of \cite[Def. 2.1.3]{DF3}.
 Recall that one can define the \emph{Euler class} $e(\cE)$ as the image of the Thom class $\thom(\cE)$
 under the canonical map:
$$
\E^{2r,r}(\Th(\cE)) \rightarrow \E^{2r,r}(\cE) \xrightarrow{\sim} \E^{2r,r}(X).
$$
Moreover, when $\E$ admits a Thom structure, one deduces an isomorphism:
\begin{equation}\label{coh_HPinfty}
\E^{n,m}(\HP^\infty_S) \simeq \colim_{n>0} \E^{n,m}(\HP^n_S)
\end{equation}
\emph{i.e.} the $\colim^{(1)}$ vanishes here.
 This follows from the symplectic projective bundle theorem for $\E$
 (see \cite[Th. 2.2.3]{DF3})

An $\MSp$-algebra structure on $\E$ (over $S$) will be a morphism $\varphi:\MSp_S\rightarrow \E$
 of ring spectra.
 We denote by $\Sigma_T^2 1_\E$ the image of the unit $1_\E \in \E^{00}(S)$ of the ring spectrum $\E$
 under the $(\PP^1)^{\wedge 2}$-suspension isomorphism $\E^{00}(S) \simeq \tilde \E^{4,2}(\HP^1_S)$.
\end{num}
\begin{thm}[Panin-Walter]\label{thm:PW_Sp-oriented}
Let $\E$ be a ring spectrum over $S$. Then the following data are in bijective
 correspondance:
\begin{enumerate}
\item the cohomology classes $b \in \tilde \E^{4,2}(\HP^\infty_S)$ such
 that $b|_{\HP^1_S}=\Sigma_T^2 1_\E$.
\item The $\Sp$-orientation on $\E$ over $S$ in the sense of \Cref{df:G-orientation}.
\item The $\MSp$-algebra structure on $\E$ over $S$: $\varphi:\MSp_S \rightarrow \E$
\end{enumerate}
Moreover, the bijections are constructed as follows:
\begin{itemize}
\item (3) $\rightarrow$ (2): for any symplectic bundle $\cU$ over a smooth $S$-scheme $X$,
 one defines $\thom^\E(\cU)=\varphi_*(\thom^\MSp(\cU))$.
\item (2) $\rightarrow$ (1): one defines: $b_n=\varphi_*(e(\cV_{2})) \in \E^{4,2}(\HP^n_S)$,
 where $\cV_{2}$ is the tautological rank $2$ symplectic bundle over $\HP^n_S$.
 Using formula \eqref{coh_HPinfty}, this defines a class $b \in \E^{4,2}(\HP^\infty_S)$.
\end{itemize}
%\jean{Shall we say something about the missing implication? Or maybe simply omit the sketch of the proof.}
%\fred{Je pense que ce n'est pas la peine: les deux points ci-dessus permettent juste d'être plus précis.
%  La preuve complète est bien donnée dans la référence ci-dessous.}
\end{thm}
For the proof of that theorem, we refer the reader to \cite[Th. 1.1]{PWcobord}.
The condition (1) being much more cost effective,
 and analogous to that of $\GL$-oriented ring spectra,
 we chose to use it as an equivalent replacement of \Cref{df:G-orientation} in the case $G=\Sp$
 (or also Definition 2.3.1 of \cite{DF3}).
\begin{df}\label{df:Sp-orientation}
A symplectic orientation of a ring spectrum $\E$ will be a class $b$
 as in point (1) of the previous proposition.
 We also say that $(\E,b)$ is a symplectically oriented ring spectrum over $S$.

A morphism of symplectically oriented ring spectra
 $\varphi:(\E,b^\E) \rightarrow (\F,b^\F)$ is a morphism of ring spectra
 such that $\varphi_*(b^\E)=b^\F$.
\end{df}
Note that by the above theorem, given such a morphism $\varphi$,
 one immediately deduces for an $\Sp$-bundle $\cU$ over a smooth $S$-scheme $X$,
 the relation
\begin{equation}\label{eq:Sp-morphisms&Thom_classes}
\varphi_*(\thom^\E(\cU))=\thom^\F(\cU).
\end{equation}

\begin{rem} The symplectic orientation theory is completely analogous to
 the classical orientation theory, both in $\AA^1$-homotopy (see \cite[Section 2.1]{Deg12}), 
 and in classical topology (see \cite[Definition 4.1.1]{Ra03}):
\begin{enumerate}
\item Note as a byproduct of the above theorem that $\MSp_S$ equipped with its canonical
 orientation $b^\MSp$ (coming from the above construction of Thom classes),
 is the universal symplectically oriented ring spectrum over $S$.
 That is, it is the initial object among symplectically oriented ring spectra over $S$.
\item From the previous theorem, a symplectically oriented ring spectrum is just
 an $\MSp$-algebra, and a morphism of such is simply a morphism of $\MSp$-algebras.
\item Let $f:T \rightarrow S$ be a morphism of schemes.
 Given an $\Sp$-oriented ring spectrum $(\E,b)$ over $S$,
 we immediately deduce an $\Sp$-orientation $f^*(b)$ on $f^*(\E)$ over $T$.
 This will be very useful to reduce arguments to the base scheme.
\end{enumerate}
\end{rem}

As a corollary of the 

\begin{num}\label{num:Sp-Euler}
One can interpret Euler classes in analogy with the $\GL$-oriented case.
 Given any scheme $S$, according to \cite[Prop. 1.16, p. 130]{MV},
 there exists a canonical map:
$$
\mathrm{H}^1(S,\Sp_2)  \rightarrow [S_+,\BSp_2^S]=[S_+,\HP^\infty_S].\footnote{According to \cite[4.1.2]{AHW},
 this map is an isomorphism whenever $S$
 is affine and ind-smooth over a Dedeking ring but we will not use that fact.}
$$
 The left hand-side can be interpreted
 as the set of isomorphisms classes of rank $2$ symplectic bundle $\PicSp(S)$.
 Note that this map is compatible with pullbacks in $S$. Besides, given a smooth $S$-scheme $X$,
 one similarly gets a canonical map:
$$
\mathrm{H}^1(X,\Sp_2) \rightarrow [X_+,\BSp_2^S]=[X_+,\HP^\infty_S].
$$
One should be aware of an essential difference with the $\GL$-case here:
 there is no group structure on $\PicSp(S)$. First, the tensor product of two symplectic bundles
 is an orthogonal bundle. Secondly, even if one considers triple tensor products,
 the rank becomes $8$. The correct algebraic structure will be explored
 through the notion of formal ternary laws.

However, given an $\Sp$-oriented spectrum $(\E,b)$, one gets a streamlined interpretation of the associated
 Euler classes in terms of $b$.
 First one can view $b$ as a map $b:\Sigma^\infty \HP^\infty_S \rightarrow \E(2)[4]$.
 Then given any smooth $S$-scheme $X$, one gets a natural map: 
\begin{align*}
\PicSp(X)=\mathrm{H}^1(X,\Sp_2) &\rightarrow [X_+,\HP^\infty_S] \xrightarrow{\Sigma^\infty}
 [\Sigma^\infty X_+,\Sigma^\infty \HP^\infty_S]^{st} \\
& \xrightarrow{b_*} [\Sigma^\infty X_+,\E(2)[4]]^{st}=\E^{4,2}(X).
\end{align*}
One easily checks that this is the \emph{Euler class map}, which to the class of a rank $2$
 symplectic bundle $\cU$ associates $e(\cU)$ whose construction was recalled in
 \ref{num:symplectic_Thom_Grass_Euler}. In this case, this coincides with the first
 Borel class $b_1(\cU)$ which will be recalled in the next paragraph.
\end{num}

\begin{rem}
In fact, one easily checks that the data of an $\Sp$-orientation
 on $\E$ over $S$ is equivalent to the data for any smooth $S$-scheme $X$ of
$$
b_1:\PicSp(X) \rightarrow \E^{4,2}(X)
$$
such that:
\begin{enumerate}
\item $b_1$ is natural with respect to the pullback on both functors
\item for the tautological symplectic bundle $\cV$ on $\HP^1_S$,
 one has $b_1(\cV)=(1,0)$ via the identification $\E^{4,2}(\HP^1_S)=\E^{0,0}(S) \oplus \E^{4,2}(S)$.
\end{enumerate}
\end{rem}

\begin{num}\label{num:Sp-proj-bdl}
As already said, a symplectically oriented ring spectrum $(\E,b)$ automatically satisfies
 the symplectic projective bundle theorem, \cite[Th. 2.2.3]{DF3}.
 In particular, putting $\E^{**}=\E^{**}(S)$ as a bigraded ring,
 one obtains an isomorphism of $\E^{**}$-algebra
$$
\E^{**} [[t]] \rightarrow \E^{**}(\HP^\infty_S), t \rightarrow b
$$
which is bihomogeneous, of bidegree $(0,0)$,
 when one decides that $t$ on the left hand-side has bidegree $(4,2)$.
 As a corollary of the previous theorem one therefore deduces:
\end{num}
\begin{cor}\label{cor:chg_Sp-orient}
Let $(\E,b)$ be a symplectically oriented ring spectrum over $S$.
 Then the symplectic orientations $b'$ of $\E$ are in one-to-one correspondence
 with the local parameters of bidegree $(4,2)$ of the power series ring $\E^{**}[[t]]$,
 \emph{i.e.} with the power series of the form
$$
\phi(t)=t+a_2t^2+a_3t^3+\hdots
$$
where $a_n \in \E^{4-4n,2-2n}$.  \\
 The bijection is given by $\phi \mapsto \phi(b)$ in
 $\E^{**}(\HP^\infty_S) \simeq \E^{**}[[b]]$.
\end{cor}
Note moreover that when one has a class $b'$ as above associated
 with a power series $\phi$, then for any rank $2$ symplectic
 bundle $\cU$ over a smooth $S$-scheme $X$, one gets:
\begin{equation}\label{eq:firstBorel&chg_orient}
b'_1(\cU)=\phi(b_1(\cU)) \in \E^{4,2}(X),
\end{equation}
where the right hand-side makes sense as $b_1(\cU)$ is nilpotent
 (see \cite{DF3}).

\subsection{Virtual symplectic Thom classes}\label{sec:virt_symplect}

According to \Cref{prop:stable_G-orientations},
 symplectically oriented ring spectra admits virtual symplectic Thom classes,
 associated with stable symplectic bundles. However, hermitian K-theory of symplectic bundles
 is finer as it involves the equivalence relation modulo metabolic spaces (see e.g. \cite{Kneb}).
 We will show in this subsection that symplectic Thom classes are compatible with
 the metabolic equivalence relation.

\begin{num}
We first give a quick recollection on the hermitian $Q$-construction and the forgetful functor, following \cite{Schlichting09}.

We start with a triple $(\mathcal{E},*,\eta)$, where $\mathcal{E}$ is an exact category, $(-)^*$ is a duality (i.e. an exact functor $\mathcal E^{\mathrm{op}}\to \mathcal{E}$) and a canonical isomorphism $\eta:1\to (-)^{**}$ satisfying an extra condition \cite[Definition 2.1]{Schlichting09}. Recall that a symmetric form is a pair $(X,\varphi)$, where $X$ is an object of $\mathcal{E}$ and $\varphi:X\to X^*$ is a morphism which is symmetric in the sense that the diagram
\[
\xymatrix{X\ar[r]^-{\varphi}\ar[d]_-{\eta} & X^*\ar@{=}[d] \\
X^{**}\ar[r]_-{\varphi*} & X^*}
\]
is commutative. A symmetric form is called \emph{symmetric space} if $\varphi$ is an isomorphism. A morphism of symmetric forms $f:(X,\varphi)\to (Y,\psi)$ is a morphism $f:X\to Y$ such that the diagram
\[
\xymatrix{X\ar[r]^-f\ar[d]_-{\varphi} & Y\ar[d]^-\psi \\
X^* & Y^*\ar[l]^/-2pt/{f^*}}
\]
is commutative. Let now $(X,\varphi)$ be a symmetric space (i.e. $\varphi$ is an isomorphism). Let $i:L\to X$ be an admissible monomorphism. The orthogonal of $L$, denoted by $L^{\perp}$ is the kernel of $X\xrightarrow{i^*\varphi} L^*$, which is also admissible. We say that $L$ is totally isotropic if the composite $i^*\varphi i$ is trivial and the associated morphism $L\to L^{\perp}$ is admissible. In the special case where $L\to L^{\perp}$ is an isomorphism, we say that $L$ is a Lagrangian.


We now form the $Q$-construction in this context, following \cite[\S 4.1]{Schlichting09}. We let $Q^h\mathcal E$ be the category whose objects are symmetric spaces $(X,\varphi)$ and whose morphisms are equivalences classes of diagrams of the form
\[
\xymatrix{X & U\ar[l]_-p\ar[r]^-i & Y}
\] 
where $p$ is an admissible epimorphism, $i$ is an admissible monomorphism, and the diagram
\[
\xymatrix{U\ar[r]^-i\ar[d]_-p & Y\ar[d]^-{i^*\psi} \\
X\ar[r]_-{p^*\varphi} & U^*}
\]
is bicartesian. Two such diagrams $(U,i,p)$ and $(U',i',p')$ are equivalent if there exists an isomorphim $g:U\to U'$ such that $i'g=i$ and $p'g=p$. The composition is obtained as in Quillen's original $Q$-construction, and there is thus an obvious functor
\[
Q^h(\mathcal E,*,\eta)\to Q\mathcal E
\]
which induces a morphism on the nerves of the relevant categories. To simplify the notation, we will now write $\mathcal E$ for $(\mathcal E,*,\eta)$.
 Following Schlichting, we denote by $\mathrm{GW}(\mathcal E)$ the homotopy fiber (with respect to a zero object in $\mathcal E$) of the
 induced map of (Kan) simplicial sets:
\[
N(Q^h\mathcal E) \to N(Q\mathcal E).
\]
As a special case of an exact category with duality, we may consider the category $H\mathcal E=\mathcal E\times \mathcal E^{\mathrm{op}}$. There is an obvious duality, defined on objects as $(A,B)^{\sharp}=(B,A)$ and natural isomorphism $1\to {}^{\sharp\sharp}$ the identity. If $(\mathcal E,*,\eta)$ is an exact category with duality, we define a functor 
\[
F:(\mathcal E,*,\eta)\to H\mathcal E
\] 
on objects by $A\mapsto (A,A^*)$ and on morphisms by $f\mapsto (f,f^*)$. There is a natural isomorphism
\[
F(A^*)=(A^*,A^{**})\to F(A)^{\sharp}=(A^*,A)
\]
given by the pair $(\mathrm{Id}_A,\eta_A)$. Consequently, we obtain a commutative diagram
\[
\xymatrix{Q^h\mathcal E\ar[r]\ar[d] &  Q\mathcal E\ar[d]\\
Q^hH\mathcal E\ar[r] &  QH\mathcal E
}
\] 
Now, there is an equivalence of categories $Q^hH\mathcal E\to Q\mathcal E$ given on objects by $(A,B)\mapsto A$ and on morphisms by $(f,g)\mapsto f$ (see \cite[Remark 4.3]{Schlichting09}).
 It follows that the homotopy fiber of 
\[
N(Q^hH\mathcal E) \to N(QH\mathcal E) \simeq N(Q\mathcal E) \times N(Q\mathcal E)
\]
is $\Omega N(Q\mathcal E):=\mathrm{K}(\mathcal E)$ (see \cite[Proposition 4.7]{Schlichting09}).
 Consequently, we obtain the \emph{forgetful map}
\[
F:\mathrm{GW}(\mathcal E)\to \mathrm{K}(\mathcal E).
\]
Our main example here is the category $\mathcal E=\mathcal{V}(X)$ of vector bundles over a scheme $X$, with usual duality $V^*=\mathrm{Hom}(V,\OO_X)$ and canonical isomorphism $\eta=-\varpi$, where $\varpi:V\to V^{**}$ is the usual identification (given by the evaluation). In the usual notation, this gives the \emph{forgetful map}
\[
F:\mathrm{KSp}(X)\to  \mathrm{K}(X).
\]
We can view $F$ as a morphism of $\infty$-categories, and consider the induced map on the associated homotopy groupoids:
$$
f_\Sp:\uKSp(X) \rightarrow \uK(X).
$$
The pseudo-functoriality of the above construction actually gives
 a stable structure of degree $2$ in the sense of \Cref{df:stable_structure}.
\end{num}
\begin{df}\label{df:fine-sp-structure}
We will call the stable structure $f_\Sp:\uKSp \rightarrow \uK$ over $\base$ defined above
 the \emph{fine symplectic stable structure}.
\end{df}

\begin{num}
Note that, by construction, one gets a commutative diagram:
$$
\xymatrix@R=-2pt@C=40pt{
\uK^{\Sp}(X)\ar^{\sigma_\Sp}[rd]\ar_{\kappa_X}[dd] & \\
& \uK(X) \\
\uKSp(X)\ar_{f_\Sp}[ru]
}
$$
where $\sigma_\Sp$ is stable structure associated with
 the symplectic group in \Cref{ex:stable_G-structure}.
 Moreover, $\kappa_X$ is compatible with pullbacks in $X$, or in other words,
 corresponds to a morphism of $\base$-fibred categories in Picard groupoid
 (see \Cref{num:fibred_cat&K}).
 The next lemma says that Zariski locally, both stable structures coincide.
\end{num}
\begin{lm}
If $X$ an affine scheme, then the functor $\kappa_X$ is an equivalence of categories.
\end{lm}
\begin{proof}
Indeed, by construction, the map induced by $\pi$ on the respective sets of isomorphism classes coincides
 with the quotient map
$$
\mathrm K^{\Sp}_0(X) \rightarrow \mathrm K^{\Sp}_0(X)/\mathcal I
$$
where $\mathcal I$ is the ideal generated by metabolic symplectic spaces.
 Then the result follows as in \cite[\textsection 4, Prop. 1]{Kneb}.
\end{proof}

We have seen in \Cref{prop:stable_G-orientations} that,
 on a given ring spectrum $\E$ over $S$,
 the data of a  symplectic orientations on $\E$ over $S$-schemes
 is equivalent to that of a $\sigma_\Sp$-orientations on $\E$.
 Thanks to the preceding result, we can make a finer statement.
\begin{prop}\label{prop:sp-or&KSp-or}
Let $\E$ be a ring spectrum over $S$.
 Then there are bijective correspondences between:
\begin{enumerate}
\item The symplectic orientations on $\E$ over $S$ in the sense of \Cref{df:Sp-orientation}.
\item The $\sigma_\Sp$-orientations on $\E$ over $S$ in the sense of \Cref{df:sigma-orientation}.
\item The $f_\Sp$-orientations on $\E$ over $S$ in the sense of \Cref{df:sigma-orientation}.
\end{enumerate}
Moreover, the map between (3) and (2) is simply obtained by composing with
 the pseudo-natural transformation $\kappa:\uK^\Sp \rightarrow \uKSp$.
\end{prop}
\begin{proof}
In view of \Cref{prop:stable_G-orientations}, we need only proving the equivalence
 between datas (2) and (3). Recall that, according to Jouanolou's trick,
 given any scheme $X$, there exists a (Zariski-local) $\AA^1$-torsor
 $p:X' \rightarrow X$
 such that $X'$ is affine.
 Then the result follows from the preceding lemma applied to $X'$,
 and the fact the pullback map: $p^*:\E\modd_X \rightarrow \E\modd_{X'}$
 is an equivalence of categories.
\end{proof}

\begin{num}\textit{Virtual symplectic Thom classes}.
Following \Cref{num:Thom_iso_explicit},
 we now make explicit the data of the $f_\Sp$-orientation
 associated with an $\Sp$-oriented ring spectrum $(\E,\thom)$ over $S$.

Given a stable symplectic bundle $\omega \in \uKSp(X)$,
 of rank $r:=\rk(f_\Sp(\omega))$, and with Thom space $\Th(\omega):=\Th(f_\Sp(\omega))$ in $\SH(X)$,
 one gets a \emph{virtual Thom class}:
$$
\thom(\omega) \in \E^{2r,r}(\Th(\omega))
$$
such that the natural map:
$$
\E^{**}(X) \rightarrow \E^{*+2r,*+r}(\Th(\omega)), \lambda \mapsto \lambda.\thom(\omega)
$$
is an isomorphism of bigraded $\E^{**}(X)$-modules.

Moreover, this class is compatible with isomorphisms in $\omega$, pullbacks in $X$
 in an obvious sense. Let us spell out explicitly the compatibility with 
 sums of two stable symplectic bundles $\omega, \omega'$ of respective ranks $r$ and $r'$.
 One gets a natural product map:
$$
\mu:\E^{2r,r}(\Th(\omega)) \otimes_\ZZ \E^{2r',r'}(\Th(\omega'))
  \rightarrow \E^{2r+2r',r+r'}(\Th(\omega+\omega')).
$$
Then the following relation holds: $\mu(\thom(\omega) \otimes \thom(\omega'))=\thom(\omega+\omega')$.
 We can equivalently write this as:
$$
\thom(\omega).\thom(\omega')=\thom(\omega+\omega').
$$
One deduces that $\mu$ induces an isomorphism of bigraded $\E^{**}(X)$-modules:
$$
\E^{**}(\Th(\omega)) \otimes_{\E^{**}(X)} \E^{*}(\Th(\omega'))
 \xrightarrow \sim \E^{**}(\Th(\omega+\omega')).
$$
Finally the virtual Thom class extends the Thom class associated with
 a symplectic vector bundle. Given a symplectic bundle $(V,\psi)$,
 one associates a virtual symplectic bundle $[V,\psi] \in \uKSp(X)$
 such that $f_\Sp([V,\psi])=[V]$ the stable vector bundles assocaited with $V$.
 One further gets an identification $\Th([V,\psi])=\Th([V])=\Th(V)$,
 by construction of the Thom space of the virtual bundle $[V]$.
 Given this identification, one has in the group $\E^{2r,r}(\Th(V))$ with $r=\rk(V)$:
$$
\thom([V,\psi])=\thom(V,\psi).
$$
\end{num}



\subsection{Higher Borel classes and $\GW$-linearity}

\begin{num}\textit{Borel classes}. \label{num:Borel_Sp}
Recall that to a symplectic orientation $b$ of $\E$ is associated by Panin and Walter
 a theory of Borel classes $b_n(\cU) \in \E^{4n,2n}(X)$, $\cU/X$ a symplectic vector bundle
 over a smooth $S$-scheme $X$. Besides, the Euler class introduced in Paragraph
 \ref{num:symplectic_Thom_Grass_Euler}
 satisfies the relation: $e(\cU)=b_r(\cU)$ when $\cU$ has rank $2r$.
 We refer the reader to \cite[Section 2.2]{DF3}.
 Following the classical scheme, one defines the total Borel class associated with
 $\cU$ as the following polynomial expression in $t$:
\begin{equation}\label{eq:total_borel}
b_t(\cU)=\sum_{n=0}^r b_n(\cU).t^r.
\end{equation}
Note moreover that, using relation \eqref{eq:Sp-morphisms&Thom_classes} and the construction
 of Borel classes, one deduces further that for a morphism $\varphi$ as in \Cref{df:Sp-orientation},
 we have the following relation in $\F^{4n,2n}(X)$:
\begin{equation}\label{eq:Sp-morphisms&Borel_classes}
\varphi_*(b_n^\E(\cU))=b_n^\F(\cU).
\end{equation}
\end{num}

\begin{num}\label{num:sp-splitting-principle}
\textit{The symplectic splitting principle}. As for Chern classes,
 one has a splitting principle for Borel classes.
 Indeed, given any symplectic bundle $\cV$ over a scheme $X$, there exists an affine morphism $p:X' \rightarrow X$
 inducing a monomorphism $p^*:\E^{**}(X) \rightarrow \E^{**}(X')$ with the property that $p^{-1}(\cV)$ splits
 as a direct sum of rank $2$ symplectic bundles: $p^{-1}(\cV)=\mathcal X_1 \perp \hdots \perp \mathcal X_n$.

One defines the {\it Borel roots} of $\cV$ as the Borel classes $\xi_i=b_1(\mathcal X_i)$ so that
 by the Whitney sum formula \cite[Theorem 10.5]{PWQuaternionic} the Borel classes of $\cV$ are the elementary symmetric polynomials
 in the variables $\xi_i$.

As in the classical case, one can compute universal formulas involving Borel classes of $\cV$ by introducing Borel roots
 $\xi_i$, which reduces to rank $2$ symplectic bundles, compute as if $\cV$ was completely split, and
 then express the resulting formula in terms of
 the elementary symmetric polynomials in the $\xi_i$. %This principle will be used repeatedly in Section \ref{sec:ternary}.
\end{num}

\begin{num}
We now give a complement to the theory of Borel classes which was known previously only under the stronger assumption of the existence of an $\SL$-orientation
(see \cite[2.2.8]{DF3}). To start with, let us recall the following definition.

\begin{df}
Let $\lambda\in \OO(S)^\times$. We denote by $\langle \lambda\rangle\in \mathrm{End}_{\mathcal{SH}(S)}(\mathbf{1})$ the endomorphism induced by the pointed morphism $\lambda:\PP^1_{S}\to \PP^1_{S}$ given by $[x:y]\mapsto [\lambda x:y]$.
\end{df}

As a result, $\lambda$ acts on the set $[X,Y]_{\mathcal{SH}(S)}$ for any objects $X,Y$ via 
\[
f\mapsto f\wedge \langle \lambda\rangle:X\wedge 1\to Y\wedge 1.
\]
By abuse of notation, we still denote by $\langle\lambda\rangle$ the endomorphism $\mathrm{Id}_X\wedge \langle\lambda\rangle$. The aim of this section is to prove the following theorem.

\begin{thm}\label{thm:epsilon}
Let $(\E,b)$ be an $\Sp$-oriented ring spectrum, and thom classes $\thom$.
 Then for any unit $\lambda \in \mathcal O_S(X)^\times$, 
 and any $\Sp$-vector bundle $(V,\psi)$ over $X$, the following relation holds:
\[
\thom(V,\lambda\psi)=\langle \lambda \rangle \thom(V,\psi).
\]
\end{thm}

As a corollary, we obtain the following result whose proof is the same as \cite[2.2.8]{DF3}.

\begin{cor}\label{cor:epsilonlinearity}
Let $(\E,b)$ be an $\Sp$-oriented ring spectrum and let $\lambda \in \mathcal O_X(X)^\times$ be a unit. Then
\[
b_i(V,\lambda\psi)=\langle \lambda^i\rangle b_i(V,\psi)
\]
for any $\Sp$-vector bundle $(V,\psi)$ over $X$ and any $i\in\NN$.
\end{cor}

The proof of the theorem will occupy the remaining of this section.
 We first reduces to the case $\E=\MSp_S$ (using \Cref{thm:PW_Sp-oriented})
 and to the case where of the tautological symplectic bundle $(U,\varphi)$ on $\mathrm{BSp}_2$
 seen over the base scheme $S$. Given $\lambda \in \cO_S(S)^\times$, the relation to prove is now:
\[
\mathrm{th}(U,\lambda\varphi)=\langle \lambda\rangle \mathrm{th}(U,\varphi).
\]
The proof of the theorem will require some elaborate computations with the hyperbolic Grassmannians. Recall that $\HP^n_S=\HGr_S(2,\omega_{2n+2})$ and that there are embeddings
\[
i_n:\HP^n_S\to \HP^{n+1}_S
\]
for any $n\in\NN$. In particular, $\HP^0_S=S$ and we will consider $\HP^{n}_S$ as pointed by the image of $\HP^0_S=S$ under the above map. In the sequel, we will suppress the subscript $S$ from the notation to slightly lighten the expressions. 

There is an interesting cell structure on $\HP^n$ for any $n\in\NN$, as explained by Panin-Walter in \cite{PWQuaternionic}. There is a decomposition
\[
\HP^n=X_0\sqcup X_2\sqcup \ldots\sqcup X_{2n}
\]
where each $X_i$ is of codimension $2i$, smooth and quasi-affine. The closure $\overline X_{2i}$ of $X_{2i}$ is of the form $\overline X_{2i}=X_{2i}\sqcup X_{2i+2}\sqcup\ldots\sqcup X_{2n}$ and there is a projection
\[
q_i:\overline X_{2i}\to \HP^{n-i}
\]
identifying $\overline X_{2i}$ with the total space of the universal bundle $\mathcal U_2$ over $\HP^{n-i}$. The composite of the zero section $\HP^{n-i}\to \overline X_{2i}$ with $\overline X_{2i}\to \HP^{n}$ is the composite of embeddings $i_{n-1}\circ\ldots\circ i_{n-i}$ described above. In more detail, write $e_1,\ldots,e_{2n+2}$ for the usual basis of $\OO_S^{2n+2}$. Observe that $E_i=\mathrm{Vect}(e_{2j-1},j=1,\ldots,i)$ is a totally isotropic subspace, in the sense that $E_i\subset E_i^\perp$, and that we have a filtration of $\OO_S^{2n+2}$ of the form
\[
0:=E_0\subset E_1\subset E_{n+1}=E_{n+1}^\perp\subset \ldots\subset E_1^\perp\subset 0^\perp=\OO_S^{2n+2}
\]
We have for any $i=0,\ldots,n$ an embedding $\mathrm{Gr}(2,E_i^{\perp})\subset \mathrm{Gr}(2,\OO_S^{2n+2})$ and we set $\overline X_{2i}:=\mathrm{Gr}(2,E_i^{\perp})\cap \HP^{n+1}$. In terms of the universal sequence, we can interpret these closed subschemes as follows. Consider the projection $p_2:\OO_S^{2n+2}\to \OO_S$ on the second factor. As $E_1^\perp=\langle e_1,e_3,e_4,\ldots,e_{2n+2}\rangle$ we see that $\overline X_2$ is the closed subscheme where the section
\[
\mathcal{U}\xrightarrow{i} \OO^{2n+2}\xrightarrow{p_2} \OO
\]
vanishes. In particular, the normal bundle to $\overline X_2$ in $\HP^n$ is $(\mathcal{U}^\vee)_{\vert \overline X_2}$. Besides, if $p_1:\OO_S^{2n+2}\to \OO_S$ is the projection to the first factor, it is easy to check that if the above composite vanishes (i.e. we are on $\overline X_2$) then the composite 
\[
\mathcal{U}\xrightarrow{i} \OO^{2n+2}\xrightarrow{p_1} \OO
\]
is arbitrary (i.e. the symplectic form on $\mathcal{U}$ is automatically nondegenerate). Observing that the restriction of $\mathcal{U}$ to $\HP^{n-1}$ is the universal bundle, this realizes $\overline X_2$ as the total space of $\mathcal{U}^\vee$ over $\HP^{n-1}$. We also deduce from this that the normal bundle to $\HP^{n-1}$ in $\HP^{n}$ is $\mathcal{U}^\vee\oplus \mathcal{U}^\vee$. To conclude, let us say that the filtration $\HP^{n-1}=X^\prime_0\sqcup\ldots \sqcup X^\prime_{2n-2}$ pulls back to the filtration $X_2\sqcup\ldots \sqcup X_{2n}$ under the projection $q_{n-1}:\overline X_2\to \HP^{n-1}$. 

Let us now concentrate on the geometry of $X_{2i}$ for $i=0,\ldots,n$, starting with $X_0$. By definition, $X_0$ is the open subscheme over which the composite 
\[
\pi:\mathcal{U}\xrightarrow{i} \OO^{2n+2}\xrightarrow{p_2} \OO
\]
is non trivial. This yields a nowhere vanishing section of $\mathcal{U}$ with kernel $\OO$ (as $\mathcal U$ has trivial determinant). We set $Y=\{f\in U\vert \pi(y)=1\}$. This is a $\ker(\pi)$-torsor over $X_0$, i.e. an $\AA^1$-torsor over $X_0$ of the form $\mu:Y\to X_0$. We now work over $Y$. It is easy to see that there exists a unique global section $e$ of $\mu^*\mathcal{U}$ which has the property that the symplectic form on $\OO^2$ induced by the isomorphism $j:\OO^2\simeq \mathcal{U}$ obtained via $f,e$ is the usual hyperbolic one. 

We may then write $i:\mathcal{U}\to \OO^{2n+2}$ as a matrix
\[
\begin{pmatrix} 1 & 0 \\ a & b \\ v & w\end{pmatrix}
\]
where $a,b\in \AA^1$ and $v,w\in \AA^{2n}$. The restriction of $\sigma_{2n+2}$ under this map is the matrix
\[
\begin{pmatrix} 0 & b+v^t\sigma_{2n}w \\ -b+w^t\sigma_{2n}v & 0\end{pmatrix}
\]
and it follows that $b+v^t\sigma_{2n}w=1$. Now, the action of $\mathbb{G}_a$ on $(a,v,w)$ is of the form $\mathbb{G}_a\times \AA^1\times \AA^{2n}\times \AA^{2n}\to \AA^1\times \AA^{2n}\times \AA^{2n}$  given by 
\[
\lambda\cdot \begin{pmatrix} 1 & 0 \\ a & b \\ v & w\end{pmatrix}=\begin{pmatrix} 1 & 0 \\ a & b \\ v & w\end{pmatrix}\begin{pmatrix} 1 & 0 \\ \lambda & 1\end{pmatrix},
\]
i.e. $\lambda\cdot (a,v,w)=(a+b\lambda,v+\lambda w,w)=(a+\lambda(1-v^t\sigma_{2n}w),v+\lambda w,w)$. We may extend this description to the other open cells $X_{2i}$ (or use inductively the fact that the pull-back of $X^\prime_0\sqcup\ldots \sqcup X^\prime_{2n-2}$ is $X_2\sqcup\ldots \sqcup X_{2n}$) to obtain the following result (\cite[Theorem 3.4]{PWQuaternionic}).

\begin{lm}
The cell $X_{2i}$ is the quotient of $\AA^1\times \AA^i\times \AA^{2n-2i}\times \AA^i\times \AA^{2n-2i}$ under the free action of $\mathbb{G}_a$ given by
\[
\lambda\cdot (a,v_1,v_2,w_1,w_2)\mapsto (a+\lambda(1-v_2^t\sigma_{2n-2i}w_2),v_1+\lambda w_1,v_2+\lambda w_2,w_1,w_2).
\]
In particular, the cells $X_{2i}$ are $\AA^1$-contractible for $i=0,\ldots,n$.
\end{lm}

As a consequence, we obtain the following lemma. 

\begin{lm}
For $n\geq 1$, there is a canonical weak-equivalence 
\[
c_n:\HP^n\simeq \mathrm{Th}(U_{\vert \HP^{n-1}}).
\]
\end{lm}

\begin{proof}
Consider the open immersion $X_0\subset \HP^n$, whose closed complement is $\overline X_2$. The normal bundle to $\overline X_2$ in $\HP^n$ being $\mathcal{U}^\vee$, we obtain by purity a cofiber sequence
\[
X_0\to\HP^n\to  \mathrm{Th}(U^\vee_{\vert\overline X_2})
\]
and therefore, as $X_0$ is contractible, a weak-equivalence $\HP^n\to  \mathrm{Th}(U^\vee_{\overline X_2})$. The inclusion $\HP^{n-1}\subset \overline X_2$ is a weak-equivalence (being the zero section of a vector bundle) and the pull back of $U^\vee_{\vert\overline X_2}$ under this map is just $U^\vee_{\HP^{n-1}}$, which is isomorphic to its dual. We then obtain a string of canonical weak-equivalences $\mathrm{Th}(U^\vee_{\vert\overline X_2})\simeq \mathrm{Th}(U^\vee_{\vert \HP^{n-1}})\simeq \mathrm{Th}(U_{\vert \HP^{n-1}})$.
\end{proof}

For $n\geq 0$, the universal subbundle $\mathcal U$ on $\HP^{n+1}$ restricts to the canonical bundle on $\HP^n$ and we obtain an induced map 
\[
j_n:\mathrm{Th}(U_{\vert \HP^{n}})\to \mathrm{Th}(U_{\vert \HP^{n+1}}).
\]
The following lemma will prove useful in the sequel. 

\begin{lm}\label{lem:fundamental}
Let $n\geq 1$ and let $\tau_n:\HP^n\to \mathrm{Th}(U_{\vert \HP^{n}})$ be the canonical morphism.
Then, the following diagram commutes in $\mathcal{H}(S)$
\[
\xymatrix{\HP^n\ar[r]^-{c_n}\ar[d]_-{i_n}\ar[rd]_-{\tau_n} & \mathrm{Th}(U_{\vert \HP^{n-1}})\ar[d]^-{j_{n-1}}  \\
\HP^{n+1}\ar[r]_-{c_{n+1}} &  \mathrm{Th}(U_{\vert \HP^{n}}).
}
\]
\end{lm}

\begin{proof}
We first prove that the bottom triangle commutes. Consider the composite of the zero section $\overline X_2\to U_{\vert \overline X_2}$ with the projection $U_{\vert \overline X_2}\to\mathrm{Th}(U_{\vert \overline X_2})$. The deformation to the normal cone shows that the following diagram commutes in $\mathcal{H}(S)$ 
\[
\xymatrix{ & \overline X_2\ar[r]\ar[d] & U_{\vert \overline X_2}\ar[d] \\
X_0\ar[r] & \HP^{n+1}\ar[r] & \mathrm{Th}(U_{\vert \overline X_2}).}
\]
On the other hand, we have a commutative diagram
\[
\xymatrix{\HP^{n}\ar[r]\ar[d] & U_{\vert \HP^n}\ar[d] \\
\overline X_2\ar[r] & U_{\vert \overline X_2}}
\] 
inducing the weak-equivalence $\mathrm{Th}(U_{\vert \HP^n})\to \mathrm{Th}(U_{\vert \overline X_2})$. The claim follows from the fact that the composite 
\[
\HP^{n}\to \overline X_2\to \HP^{n+1}
\]
is $i_n$. We prove next that the other triangle commutes. For this, recall that $X_0^\prime\subset \HP^{n}$ is defined as the complement of the vanishing locus of a section $s:\mathcal{U}\to \OO$. We then obtain a commutative diagram
\[
\xymatrix{X_0^\prime\ar[r]\ar[d] & \HP^{n}\ar[r]\ar[d]_-s & \mathrm{Th}(U_{\vert \overline X_2^\prime})\ar@{-->}[d] \\
U\setminus 0\ar[r] & U\ar[r] & \mathrm{Th}(U).
}
\]
In $\mathcal{H}(S)$, $s$ and the zero section are equal, and we conclude observing that the composite 
\[
\xymatrix{\mathrm{Th}(U_{\vert \HP^{n-1}})\ar[r] & \mathrm{Th}(U_{\vert \overline X_2^\prime})\ar@{-->}[r] & \mathrm{Th}(U)}
\]
is just $j_n$.
\end{proof}

We now pass to the definition of an action of $\mathbb{G}_m$ on $\HP^n$ for any $n\in\NN$. Let $\lambda\in \OO(S)^\times$ and let $m_\lambda:\OO_S^2\to \OO_S^2$ be given by the matrix
\[
\begin{pmatrix} \lambda & 0 \\ 0 & 1\end{pmatrix}.
\]
A direct computation shows that $m_\lambda^\vee \sigma_2 m_\lambda=\lambda\cdot \sigma_2$ and it follows that $m_{\lambda}$ induces a map
\[
\mathrm{Sp}_2(S)\to \mathrm{Sp}_2(S)
\]
given by $A=\begin{pmatrix} a & b \\ c & d\end{pmatrix}\mapsto m_{\lambda}^{-1}\cdot A\cdot m_{\lambda}=\begin{pmatrix} a & \lambda^{-1}b \\ \lambda c & d \end{pmatrix}$.
This induces an action of $\mathbb{G}_m$ on $\mathrm{Sp}_2$, and consequently an action of $\mathbb{G}_m$ on $\mathrm{BSp}_2$. This could be defined at the level of $\HP^n$ for $n\in\NN$ as follows: Set
\[
\mathrm{Gr}(2,\OO_S^{2n+2})\to \mathrm{Gr}(2,\OO_S^{2n+2})
\]
by considering the composite 
\[
\mathcal{U}\xrightarrow{i} \OO^{2n+2}\xrightarrow{m_{\lambda}^{\oplus n+1}} \OO^{2n+2}.
\]
If $i^\vee \sigma_{2n+2} i$ is non degenerate, then so is the composite 
\[
i^\vee (m_{\lambda}^{\oplus n+1})^\vee \sigma_{2n+2} (m_{\lambda}^{\oplus n+1}i)=\lambda\cdot  (i^\vee \sigma_{2n+2} i).
\]
Consequently, the above action of $\mathbb{G}_m$ on $\mathrm{Gr}(2,\OO_S^{2n+2})$ restricts to an action of $\mathbb{G}_m$ on $\mathrm{HGr}(2,\OO_S^{2n+2})=\HP^{n}$. Let us denote by 
\[
\mu:\mathbb{G}_m\times \HP^{n} \to \HP^{n}
\]
this action, and $\pi:\mathbb{G}_m\times \HP^{n}\to \HP^{n}$ the projection to the second factor.

\begin{lm}
There is a canonical isomorphism $f:\pi^*\mathcal{U}\to \mu^*\mathcal{U}$ on $\mathbb{G}_m\times \HP^{n}$.
\end{lm}

\begin{proof}
We start by exhibiting a $\mathbb{G}_m$-equivariant open covering $\cup_{i=1}^m V_i$ of $ \HP^{n}$ which has the following properties
\begin{enumerate}
\item $V_i$ is affine for any $i=1,\ldots,m$.
\item The restriction of $\mathcal U$ to $V_i$ is free for any $i=1,\ldots,m$.
\end{enumerate} 
To obtain this, we consider the composite morphism
\[
\HP^n\to \mathrm{Gr}(2,\OO_S^{2n+2})\to \mathbb{P}_S^{\frac{(2n+2)(2n+1)}2}
\]
where the second map is the Pl\"ucker embedding. We note that the $\mathbb{G}_m$-action on $\HP^n$ is the restriction of a $\mathbb{G}_m$-action on $\mathrm{Gr}(2,\OO_S^{2n+2})$. The latter can be extended (taking exterior powers) to an action on $\mathbb{P}_S^{\frac{(2n+2)(2n+1)}2}$ which is easily seen to be diagonal. It follows that the usual covering of the latter by affine spaces (complement of the zero locus of a given coordinate) is $\mathbb{G}_m$-equivariant. We define the open subschemes $V_i\subset \HP^n$ as the intersections of this covering with $\HP^n$. The latter being affine, it follows that (1) above is satisfied. For (2), we observe that the condition on the minors is equivalent to the composite
\[
\mathcal{U}\xrightarrow{i}\OO^{2n+2}\to \OO^2
\]
being an isomorphism (the projection corresponds to the identification of some minor). 

This covering being obtained, we now choose over each $V_i$ a symplectic basis of $\mathcal U_{\vert V_i}$ and we may recover $\mathcal U$ from the open covering $V_i$ and transition functions $V_i\cap V_j\to \mathrm{Sp}_2$. A direct computations shows that $\mu^*\mathcal U$ is then obtained via the transition functions
\[
V_i\cap V_j\to \mathrm{Sp}_2\to \mathrm{Sp}_2
\]
where the second map is given by 
\[
\begin{pmatrix} a & b \\ c & d\end{pmatrix}\mapsto \begin{pmatrix} a & tb \\ t^{-1}c & d\end{pmatrix}=\begin{pmatrix} t & 0 \\ 0 & 1\end{pmatrix}\begin{pmatrix} a & b \\ c & d\end{pmatrix}\begin{pmatrix} t^{-1} & 0 \\ 0 & 1\end{pmatrix}
\]
where $t$ is a parameter of $\mathbb{G}_m$. We now define $\pi^*\mathcal U_{\vert V_i}\to \mu^*\mathcal U_{\vert V_i}$ using the matrix $\begin{pmatrix} t & 0 \\ 0 & 1\end{pmatrix}$. A direct computation shows that this behaves well with the transition functions, yielding the required isomorphism $f:\pi^*\mathcal{U}\to \mu^*\mathcal{U}$.
\end{proof}

We now fix $\lambda\in \OO(S)^\times$, obtaining a multiplication by $\lambda$ morphism
\[
\mu_{\lambda}:\HP^n\to \HP^n.
\]
Using the above lemma, we obtain a commutative diagram
\[
\xymatrix{U\ar[r]^-f\ar[rd] & \mu_{\lambda}^*U\ar[r]\ar[d] & U\ar[d] \\
 & \HP^n\ar[r]_-{\mu_{\lambda}} & \HP^n}
\]
and then a morphism of Thom spaces $\mu_{\lambda}^f:\mathrm{Th}(U_{\vert \HP^n})\to \mathrm{Th}(U_{\vert \HP^n})$ making the diagram 
\[
\xymatrix{\HP^n\ar[r]^-{\mu_{\lambda}}\ar[d] & \HP^n\ar[d] \\
\mathrm{Th}(U_{\vert \HP^n})\ar[r]_-{\mu_{\lambda}^f} & \mathrm{Th}(U_{\vert \HP^n})}
\] 
commutative. 

We note here that the action of $\mathbb{G}_m$ on $\HP^n$ stabilizes the cellular filtration  
\[
\HP^n=X_0\sqcup \ldots\sqcup X_{2n}
\] 
in the sense that the subschemes $X_i$ (and their closure) are $\mathbb{G}_m$-equivariant. As a result, we see that the morphisms in Lemma \ref{lem:fundamental} all commute with the relevant morphisms $\mu_{\lambda}$ and $\mu_{\lambda}^f$.

\begin{lm}\label{lem:initialize}
The endomorphism $\mu_{\lambda}:\HP^1\to \HP^1$ is equal to $\langle\lambda\rangle$ in $\mathrm{End}_{\mathcal{SH}(S)}(\HP^1)$.
\end{lm}

\begin{proof}
By the above remark, we have a commutative diagram
\[
\xymatrix{\HP^1\ar[r]^-{c_1}\ar[d]_-{\mu_{\lambda}} & \mathrm{Th}(U_{\vert \HP^0})\ar[d]^-{\mu_{\lambda}^f} \\
\HP^1\ar[r]_-{c_1} & \mathrm{Th}(U_{\vert \HP^0})}
\]
Now, $U_{\vert \HP^0}$ is trivial and $f$ is just the multiplication by $\begin{pmatrix} \lambda & 0 \\ 0 & 1\end{pmatrix}$. The claim follows easily. 
\end{proof}

By construction, there is a tautological Thom class 
\[
\mathrm{th}(U_{\vert \HP^n}):\mathrm{Th}(U_{\vert \HP^n})\to \mathbf{MSp}(2)[4]
\]
which has the property that the following diagram commutes for any $n\in\NN$:
\[
\xymatrix@C=5em{\mathrm{Th}(U_{\vert \HP^n})\ar[r]^-{\mathrm{th}(U_{\vert \HP^n})}\ar[d]_-{j_n} & \mathbf{MSp}(2)[4]\ar@{=}[d] \\
\mathrm{Th}(U_{\vert \HP^{n+1}})\ar[r]_-{\mathrm{th}(U_{\vert \HP^{n+1}})} & \mathbf{MSp}(2)[4]}
\]

\begin{prop}
For any $n\geq 0$, the composite
\[
\mathrm{Th}(U_{\vert \HP^n})\xrightarrow{\mu_{\lambda}^f} \mathrm{Th}(U_{\vert \HP^n})\xrightarrow{\mathrm{th}(U_{\vert \HP^n})} \mathbf{MSp}(2)[4]
\]
is equal to $\langle \lambda\rangle \mathrm{th}(U_{\vert \HP^n})$ in $[\mathrm{Th}(U_{\vert \HP^n}),\mathbf{MSp}(2)[4]]_{\mathcal{SH}(R)}$.
\end{prop}

\begin{proof}
We work by induction on $n$. If $n=0$, this is a direct consequence of (the proof of) Lemma \ref{lem:initialize}. Let us then suppose that the result is true for some $n-1\geq 0$ and consider the composite
 \[
\mathrm{Th}(U_{\vert \HP^n})\xrightarrow{\mu_{\lambda}^f} \mathrm{Th}(U_{\vert \HP^n})\xrightarrow{\mathrm{th}(U_{\vert \HP^n})} \mathbf{MSp}(2)[4].
\]
We have a commutative diagram
\[
\xymatrix{\mathrm{Th}(U_{\vert \HP^{n-1}})\ar[r]^-{\mu_{\lambda}^f}\ar[d]_-{j_{n-1}} & \mathrm{Th}(U_{\vert \HP^{n-1}})\ar[r]^-{\mathrm{th}(U_{\vert \HP^{n-1}})}\ar[d]_-{j_{n-1}} & \mathbf{MSp}(2)[4]\ar@{=}[d]
\\ \mathrm{Th}(U_{\vert \HP^n})\ar[r]_-{\mu_{\lambda}^f} & \mathrm{Th}(U_{\vert \HP^n})\ar[r]_-{\mathrm{th}(U_{\vert \HP^n})} & \mathbf{MSp}(2)[4]}
\]
By induction, we know that  
\[
(\mathrm{th}(U_{\vert \HP^n})\circ \mu_{\lambda}^f)\circ j_{n-1}=(\langle \lambda\rangle \mathrm{th}(U_{\vert \HP^n}))\circ j_{n-1}
\]
and it suffices to prove that 
\[
j_{n-1}^*:[\mathrm{Th}(U_{\vert \HP^n}),\mathbf{MSp}(2)[4]]\to [\mathrm{Th}(U_{\vert \HP^{n-1}}),\mathbf{MSp}(2)[4]]
\]
is injective for $n\geq 1$. Now, we may use Lemma \ref{lem:fundamental} to obtain
\[
\tau_n^*=c_n^*\circ j_{n-1}^*
\]
and it suffices to show that $\tau_n^*$ is injective. However, it follows from \cite{PWQuaternionic} that the cofiber sequence
\[
U\setminus 0\to \HP^n \xrightarrow{\tau_n} \mathrm{Th}(U_{\vert  \HP^n})
\]
induces a short split exact sequence
\[
0\to [\mathrm{Th}(U_{\vert  \HP^n}),\mathbf{MSp}(2)[4]]\xrightarrow{\tau_n^*}[\HP^n,\mathbf{MSp}(2)[4]]\to  [U\setminus 0,\mathbf{MSp}(2)[4]]\to 0
\]
and in particular $\tau_n^*$ is injective.
\end{proof}

\begin{proof}[Proof of \Cref{thm:epsilon}]
Taking colimits, it is sufficient to prove the result for the Thom classes of $\mathcal U_{\vert \HP^n}$ for any $n\in \NN$. It follows from the above proposition that we have
\[
\langle \lambda\rangle \mathrm{th}(U_{\vert \HP^n})=\mathrm{th}(U_{\vert \HP^n})\circ \mu_{\lambda}^f.
\]
Now, the property of Thom classes reads as $\mathrm{th}(U_{\vert \HP^n})\circ \mu_{\lambda}^f=\mathrm{th}(f^*\mu_{\lambda}^*U_{\vert \HP^n})$ and a direct computation shows that $f^*\mu_{\lambda}^*(U_{\vert \HP^n},\varphi)=(U,\langle\lambda\rangle \varphi)$.
\end{proof}
 

\end{num}
%
% Given the previous remarks, we consider the case of $\MSp_S$,
% a symplectic bundle $\cE=(E,\psi)$ over $S$ and a unit $u \in \GG(S)$.
% We will compare the Borel classes of a symplectic bundle $\cE=(E,\psi)$ over $S$
% and its "twisted" version,
% $(E,u.\psi)$.
% First, we remark that we have an isometry $\omega_2\simeq u\omega_2$ given by the matrix
%\[
%\sigma_u:\begin{pmatrix} u & 0 \\ 0 & 1\end{pmatrix}.
%\]
%More precisely,
%\[
%\begin{pmatrix} u & 0 \\ 0 & 1\end{pmatrix}\begin{pmatrix} 0 & 1 \\ -1 & 0\end{pmatrix} \begin{pmatrix} u & 0 \\ 0 & 1\end{pmatrix} =\begin{pmatrix} 0 & u \\ -u & 0\end{pmatrix}.
%\]
%This yields an action of $\GG$ on $\HP^n_S$ for any $n>0$,
% sending the universal bundle $(U_{2n},\phi)$ to $(U_{2n},u\phi)$.
% Indeed, one can easily check that the symplectic form induced on $U_{2n}$ by the composite
%\[
%U_{2n}\to \AA^{2n}_S\xrightarrow{\oplus_n \sigma_u} \AA^{2n}_S
%\]
%is $u.\psi$. In the sequel, we write $\sigma_{u,n}$ in place of $\oplus_n \sigma_u$. We then obtain a morphism
%\[
%m:\GG\times \HP^n_S\to \HP^n_S
%\]
%under which $m^{-1}(U_2,\psi)=(\sigma_{u,n}(U_2),u.\psi)$.
% One easily checks that this action of $\GG$ is compatible with the inclusions $\HP^n_S\to \HP^{n+1}_S$, yielding an action of $\GG$ on $\BSp_2^S$ of the form
%\[
%m:\GG\times\BSp_2^S\to \BSp_2^S
%\]
%We let $b=b^\MSp$ be the canonical $\Sp$-orientation on $\MSp_S$. Then one obtains a class:
%\begin{align*}
%b_1(U_2,u.\psi)=b_1(m^{-1}(U_2,\psi))&=m^*b_1(U_2,\psi) \in \MSp^{2,1}(\BSp_2^S) \\
%&=\MSp^{2,1}(\HP^\infty) \simeq \MSp^{2,1}(\MSp_S).b_1(U_2)
%\end{align*}
%where the last isomorphism follows from the symplectic projective bundle theorem.
% In particular, there exists a unique class $\lambda_u \in \MSp^{2,1}(\MSp_S)$ such that: $b_1(U_2,u.\psi)=\lambda_u.b_1(U_2)$.
%
%We are then reduced to show that the class of $\lambda_u$ is multiplication by $\langle u\rangle$. In fact, we only need to prove that the multiplication by a unit defined above
%\[
%m_u:\BSp_2\to \BSp_2
%\]
%is the multiplication by $\langle u\rangle$. \fred{This assertion should be made more precise, as by definition, $<u>$ is an endomorphism
% of either the pointed sheaf $\GGx X$, the Thom space $\Th(\AA^1_X)$, of the $0$-th sphere spectrum.}
%The action defined above is compatible with the stabilization maps and we obtain a commutative diagram
%\[
%\xymatrix{\BSp_2\ar[r]^-{m_u}\ar[d] & \BSp_2\ar[d] \\
%\BSp\ar[r]_-{m_u} & \BSp}
%\]
%The map $m_u:\BSp\to \BSp$ is precisely multiplication by $\langle u\rangle$.
% Now the composite $\BSp_2\to \BSp\xrightarrow{m_u}\BSp$ can be seen as an element of $[\BSp_2,\BSp]$ (unstable here). We have 
%\[
%[\BSp_2,\BSp]=[\Sigma_{T}^{\infty}\BSp_2,\mathbf{BSp}]=[\Sigma_{T}^{\infty}\BSp_2,\mathbf{BO}(2)[4]]=\mathbf{BO}^{4,2}(\Sigma_{T}^{\infty}\BSp_2)
%\]
%The latter is just $(\mathbf{BO}^{*,*}[[b_1]])^{(4,2)}$ (the subgroup of elements of degree $(4,2)$), and we see that the class of the composite is just $\langle u\rangle b_1$. The result follows from the commutativity of the previous diagram. We have proved the following crucial proposition.
%\fred{I am a bit uncomfortable with the last step of this proof. It seems strange that the assertion reduces to the case of hermitian (or even othgonal) K-theory as there is no map
% from $\KSp$ to $\MSp$. The global plan is ok however. I believe one can di this proof directly with $\MSp$, using formula \eqref{eq:MSp_as_colim} and going back
% to the action of $\GG$ on $\Sp$. To be done...}
%\end{num}
%
%\jean{Here is a new approach:}
%
%For any $n\in\NN$, let $U_{n}$ be the universal bundle on $\HP^n_S$ which is of rank $2$. We have a universal split sequence
%\[
%0\to U_n\xrightarrow{i_n} \OO_S^{2n+2}\xrightarrow{p_n} V_{2n}\to 0
%\]
%in which $V_{2n}$ is of rank $2n$. If $\chi_n:U_n\to U_n^\vee$ is the universal form on $U_n$, we have the retraction $r_n:\OO_S^{2n+2}\to U_n$
%defined by the composite
%\[
%\OO_S^{2n+2}\xrightarrow{h_{2n+2}}(\OO_S^{2n+2})^\vee\xrightarrow{i_n^\vee}U_n^\vee\xrightarrow{\chi_n^{-1}} U_n
%\]
%and a morphism $s_n:V_{2n}^{\vee}\to \OO_S^{2n+2}$ defined by 
%\[
%V_{2n}^{\vee}\xrightarrow{p_n^\vee}(\OO_S^{2n+2})^\vee\xrightarrow{h_{2n+2}}\OO_S^{2n+2}
%\]
%such that $\psi_n:p_ns_n:V_{2n}^{\vee}\to V_{2n}$ is a skew symmetric isomorphism.
%We obtain an isomorphism
%\[
%\OO_S^{2n+2}\xrightarrow{(r_n,p_n)}U_n\oplus V_{2n}.
%\]
%For any unit $u$, we obtain a morphism (as $\mathbb{G}_m$ acts on $\HP^n_S$)
%\[
%m_u:\HP^n_S\to \HP^n_S
%\]
%inducing a morphism $m_u:U_n\to U_n$ and in turn a morphism of Thom spaces $m_u:Th(U_n)\to Th(U_n)$. The identity on $V_{2n}$ induces the identity on Thom spaces and we consider 
%\[
%Th(\OO_S^{2n+2})\simeq Th(U_n)\wedge Th(V_{2n})\xrightarrow{m_u\wedge \mathrm{Id}} Th(U_n)\wedge Th(V_{2n})\simeq Th(\OO_S^{2n+2})
%\]
%wich is induced by an isomorphism of $\OO_S^{2n+2}$. I think that its class in 
%\[
%[Th(\OO_S^{2n+2}),Th(\OO_S^{2n+2})]=[Th(\OO_S),Th(\OO_S)]
%\] 
%is the determinant (this is certainly true over $\ZZ$ and then the result holds for $u=-1$), which should be $u$. It follows that the class of $m_u$ in 
%\[
%[Th(U_n),Th(U_n)]
%\]
%is $\langle u\rangle$. One can then check that it stabilizes with respect to $n$ and thus that the induced map on Thom spaces of the universal bundles $U$ on $\BSp_2$ is also $\langle u\rangle$. Now, we consider the composite 
%\[
%\Sigma^{\infty}\BSp_2\to \Sigma^{\infty}Th(U)\xrightarrow{m_u}\Sigma^{\infty}Th(U)\to \MSp(1)[2]
%\]
%which gives the Euler class of $e(U,u\varphi)$. The commutative diagram
%\[
%\xymatrix{\Sigma^{\infty}\BSp_2\ar[r]\ar[d]^-{m_u} & \Sigma^{\infty}Th(U)\ar[d]_-{m_u} \\
%\Sigma^{\infty}\BSp_2\ar[r] & \Sigma^{\infty}Th(U)}
%\]
%allows then to conclude that $e(U,u\varphi)=\langle u\rangle e(U,\varphi)$.
%
%%If we consider instead of $i_n$ the composite
%%\[
%%U_n\xrightarrow{i_n}\OO_S^{2n+2}\xrightarrow{\sigma_{u,n}}\OO_S^{2n+2}
%%\] 
%%the above maps change accordingly and the induced map on the Thom space of $U_{n}$ is the one we want to compute. Taking the identity on the Thom space of $V_{2n}$, this is tantamount to computing the map on the Thom space of $\OO_S^{2n+2}$. The map on $U_n\oplus V_{2n}$ is given by the composite
%%\[
%%U_n\oplus V_{2n}\xrightarrow{(r_n(u),p_n)^{-1}}\OO_S^{2n+2}\xrightarrow{(r_n,p_n)}U_n\oplus V_{2n}
%%\]
%%where $r_n(u)$ is obtained as above using $\sigma_{u,n}\circ i_n$ instead of $i_n$.
%
%\begin{prop}\label{prop:GW-linearity_Borel}
%Let $(\E,b)$ be an $\Sp$-oriented spectrum over $S$.
% Let $X$ be a smooth $S$-scheme with a global unit $u \in \GG(X)$, and $(E,\psi)$ a rank $2n$ symplectic bundle over $X$.
% Then one has:
%$$
%b_1(E,u.\psi)=\langle u^n \rangle.b_n(E,\psi).
%$$
%\end{prop}
%\begin{proof}
%By universality of the $\Sp$-oriented spectrum $\MSp$, one can assume $\E=\MSp_S$.
% Using the symplectic projective bundle theorem as in \cite[2.2.8]{DF3}, one reduces to the case $X=\HP^\infty_S=\BSp_2^S$.
% \fred{Finish as suggested above...}
%\end{proof}

\subsection{The FTL associated to an $\Sp$-oriented spectrum}

\begin{num}\label{num:SP-orient->FTL}
Let $(\E,b)$ be an $\Sp$-oriented ring spectrum over $S$.
 Its ring of coefficients $\E^{**}:=\E^{**}(S)$ is a bigraded ring.
 Moreover, the action of $\epsilon=-\langle -1 \rangle$ on $\SH(S)$
 allows to equip it with a bigraded $\ZZe$-algebra structure.
 When dealing with symplectic orientations, and the forthcoming associated
 FTL, it is more natural to consider the subring
 $\E^{\tw 2*}:=(\E^{4n,2n}(S), n \in \ZZ)$,
 which is now a graded $\ZZe$-algebra.\footnote{Indeed, Borel classes live
 in degree $\tw 2*$.}

In \cite[Def. 2.3.2]{DF3}, we proved that one can associate to $(\E,b)$
 a $(3,4)$-series that we will denote by $F_t^b(x,y,z)$
 with coefficients in $\E^{\tw 2*}$.
 The method is directly imported from the $\GL$-oriented case
 (compare \cite[2.1.19]{Deg12}):
 thanks to the symplectic projective bundle theorem (\cite[2.2.3]{DF3}),
 one gets a canonical isomorphism
$$
\E^{\tw 2*}(\HP^\infty_S \times_S \HP^\infty_S \times_S \HP^\infty_S)
 \simeq \E^{\tw 2*}(S)[[x,y,z]]
$$
where $x$ (resp. $y$, $z$) are formal variables which corresponds
 to the Borel class of the tautological symplectic bundle $\cU_1$ (resp. $\cU_2$, $\cU_3$) on $\HP^\infty_S$
  over the first (resp. second, third) coordinates.
 The triple tensor product $\cU_1 \otimes \cU_2 \otimes \cU_3$ is naturally a rank $8$ symplectic bundle. 
 It follows that one can write the total Borel class \eqref{eq:total_borel}:
\begin{equation}\label{eq:Borel&FTL}
F_t^b(x,y,z):=b_t(\cU_1 \otimes \cU_2 \otimes \cU_3)
\end{equation}
as a polynomial in $t$ of degree $4$, with coefficients in $\E^{**}[[x,y,z]]$.
 In other words, the expression $F_t^b(x,y,z)$
 defines a $(3,4)$-series in the formal variables $x,y,z$ as required.
 It satisfies the symmetry and associativity axioms of \Cref{df:FTL}
 according to the analogous properties of the tensor product
 of symplectic bundles.

In \emph{op. cit.}, the Neutral and Semi-neutral element axioms were only proved
 under the additional assumption that $\E$ is $\SL$-oriented.
 Thanks to \Cref{cor:epsilonlinearity}, we can now avoid the latter assumption on $\E$.
 In fact, the mentioned corollary readily implies that the FTL $F_t^b(x,y,z)$
 satisfies the $\epsilon$-linearity axiom.
 Once this property is known, one deduces the Neutral and Semi-neutral element
 properties as in the proof of \cite[Prop. 2.3.4]{DF3}.

The above procedure allows us to deduce the following theorem
 (which enhances \cite[\textsection 2.3]{DF3} thanks to the $\epsilon$-linearity axiom).
\end{num}
\begin{prop}
The association $(\E,b) \mapsto \big(\E^{\tw 2*},F_t^b(x,y,z)\big)$ described above
 defines a canonical pseudo-functor from the scheme-fibred category of $\Sp$-oriented ring spectra
 to the ring-fibred category of formal ternary laws $\FTL$.
\end{prop}
\begin{proof}
The association on objects was described in the paragraph preceding the statement.
 The functoriality assertion follows from Formulas \eqref{eq:Borel&FTL},
 \eqref{eq:Sp-morphisms&Borel_classes}
 and the symplectic splitting principle \cite[Par. 2.2.5]{DF3}.
\end{proof}

\begin{rem}\label{rem:repFTL&BC}
To be more explicit about the pseudo-functoriality assertion (see also the forthcoming proposition),
 we mean in particular that for any morphism
 of scheme $f:T \rightarrow S$ and any s$\Sp$-oriented ring spectrum $(\E,b)$ over $S$,
 if we let $\E_T=f^*(\E)$ with the $\Sp$-orientation $b_T$ obtained by pullback
 (from the $\MSp$-algebra structure), then one gets a canonical morphism of FTL:
$$
\big(\E^{\tw 2*}(S),F_t^b\big) \rightarrow \big(\E^{\tw 2*}(T),F_t^{b_T}\big)
$$
which can be decomposed as an extension of scalars along
 the ring extension $\E^{\tw 2*}(T)/\E^{\tw 2*}(S)$ and then an isomorphism.
\end{rem}

\begin{df}
Let $S$ be a scheme. We say that a formal ternary law $(R,F_t)$ is \emph{representable over $S$}
 if there exists an $\Sp$-oriented ring spectrum $(\E,b)$ over $S$ such that $(R,F_t) \simeq (\E^{**}(S),F_t^b)$.  
 We will say that $(R,F_t)$ is \emph{weakly representable over $S$} if there exists an $\Sp$-oriented ring spectrum $(\mathbb{E},b)$
 and a (graded) ring homomorphism $\varphi: \mathbb{E}^{*,*}(S)\to R$ such that $(R,F_t)=(\E^{**},\varphi_*F_t^b)$.

Representable (resp. weakly) will mean representable (resp. weakly) over some scheme.
\end{df}

\begin{ex}\label{ex:representableFTL}
Recall \Cref{df:addmulFTL} for additive and multiplicative FTL.
\begin{enumerate}
\item The additive FTL with coefficients in $\ZZe$ (resp. $\QQe$) is representable
 by Milnor-Witt motivic cohomology spectrum $\HMW \ZZ$ over $\QQ$ (resp. $\ZZ$)
 as proved in \cite[Th. 3.3.2]{DF3}) (resp. \cite[Cor. 3.3.5]{DF3}).
 More generally, over any $\ZZ[1/6]$-scheme (any scheme) $S$, the $\Sp$-oriented ring spectrum $\HMW\ZZ_S$
 (resp. $\HMW\QQ_S$) admits the additive formal group law (see \emph{loc. cit.} and use base change
 \Cref{rem:repFTL&BC}).
\item The multiplicative FTL with coefficients in $\ZZe^{mul}$
 is representable by the (higher) Grothendieck-Witt (aka hermitian K-theory) $\GWsp_\RR$ ring spectrum
 over $\RR$, according to \cite[Th. 6.6]{FH21} (see also \cite[\textsection 3.2]{CDFH}).
 More generally, using again base change, over any $\ZZ[1/2]$-scheme $S$,
 the $\Sp$-oriented ring spectrum $\GWsp_S$ admits the multiplicative formal group law.
\end{enumerate}
\end{ex}

\begin{rem}
There is an obvious obstruction for an FTL to be representable:
 its ring of coefficients must be isomorphic to the $\tw2*$-graded part of the
 ring of coefficients of some (motivic) ring spectrum. In particular, it must be graded!
 Note on the other hand that the grading on the ring of coefficients
 of a representable FTL is compatible with the notion of degree 
 introduced in \Cref{df:FTL}.

Apart from the obvious restriction on grading,
 the problem of determining which FTLs are representable is a mystery.
 Note that it is already an open question to determine which formal group laws
 are representable by a classical complex-oriented ring spectrum, or
 even a $\GL$-oriented (motivic) ring spectrum.
\end{rem}

\begin{prop}\label{prop:sp-orient&FTL}
Let $(\E,b)$ be an $\Sp$-oriented ring spectrum with associated formal ternary law $F_t^b$.

Then there exists a one-to-one correspondence between the strict isomorphisms of $FTL$ (\Cref{df:morph_FTL})
 of the form $(Id,\phi):(\E^{\tw 2*},F_t^b) \rightarrow (\E^{\tw 2*},F_t)$
 and the $\Sp$-orientations $b'$ of $\E$.
\end{prop}
\begin{proof}
Indeed, according to \Cref{cor:chg_Sp-orient},
 the $\Sp$-orientation $b'$ corresponds to the power series $\phi$
 with coefficients in $\E^{**}$ such that $b'=\phi(b)$.
 The fact that $(Id,\phi)$ is an isomorphism from $F_t^b$ to $F_t^{b'}$
 follows from the symplectic splitting principle and Formula
 \eqref{eq:firstBorel&chg_orient}.

Given now a strict isomorphism of the form
 $(Id,\phi):(\E^{**},F_t^b) \rightarrow (\E^{**},F_t)$,
 we get that $b'=\phi(b)$ is an $\Sp$-orientation and the isomorphism $\phi$
 gives an identification between $F_t$ and $F_t^b$.
\end{proof}

As in the classical case, one gets the following notable corollary.
\begin{cor}
Given an $\Sp$-orientable ring spectrum $\E$,
 there exists a unique formal ternary law $F_t^\E$ up to strict isomorphism associated with an $\Sp$-orientation on $\E$.

Moreover, given any ring spectrum $\F$, if $\F$ admits an $\E$-algebra structure,
 then $\F$ is $\Sp$-orientable and the FTL associated with any $\Sp$-orientation on $\F$
 is uniquely strictly isomorphism to $F_t^\E$
 (after extension of scalars along $\F^{\tw 2*}/\E^{\tw 2*}$).
\end{cor}

We close this section with the following theorem.
\begin{thm}\label{thm:log_FTL}
 Any representable formal ternary law $F_t$
 with coefficients in a $\QQ$-algebra is uniquely strictly isomorphic to the additive FTL.
\end{thm}
In other words, there exists a unique strict isomorphism of FTLs
$$
\mathrm{log}_{F_t}:(R,F_t) \rightarrow \big(R,F_t^{add}\big)
$$
called the \emph{logarithm} of $F_t$.
\begin{proof}
Let $(\E,b)$ be a rational $\Sp$-oriented ring spectrum over some scheme $S$
 whose associated FTL is $F_t$.
 The rationalization $L_\QQ\E$ still represents $F_t$,
 so we can assume that $E$ is rational.
 According to \cite[Cor. 6.2]{DFJK}, the rational sphere spectrum over $S$ is isomorphic
 to the rational Milnor-Witt motivic cohomology spectrum $\HMW\QQ_S$.
 In particular $\E$ admits a canonical $\HMW\QQ_S$-algebra structure.
 Then the theorem follows from \Cref{ex:representableFTL}(1) and the preceding corollary.
\end{proof}

\subsection{Todd classes and FTL}

\begin{num}
As in \Cref{num:todd},
 we consider $\Sp$-oriented ring spectra $(\E,b)$ and $(\F,b')$
 and a morphism  of ring spectra $\psi:\E \rightarrow \F$.
 For any scheme $S$, we consider the induced morphism of bigraded rings:
$$
\psi_{\HP^\infty_S}:\E^{**}(\HP^\infty_S) \rightarrow \F^{**}(\HP^\infty_S).
$$
According to the symplectic projective bundle theorem (\cite[Th. 2.2.3]{DF3}), $\psi_{\HP^\infty_S}$ corresponds to a morphism of bigraded rings:
$$
\E^{**}(S)[[u]] \rightarrow \F^{**}(S)[[t]]
$$
where $u$ and $t$ are indeterminates of bidegree $(4,2)$.
 Let us denote by $\phi(t)$ the power series image of $u$ under this latter map.
 By definition, $\phi(t)$ is uniquely determined by the relation in $\F^{**}(\HP^\infty_S)$:
\begin{equation}
\psi_{\HP^\infty_S}(b)=\phi(b').
\end{equation}
Moreover, $\phi(b')$ is an $\Sp$-orientation of $\F$ over $S$
 and it follows from \Cref{cor:chg_Sp-orient} that it has the form:
$$
\phi(t)=t+\sum_{i>1} a_i.t^i
$$
where $a_i \in \F^{4-4i,2-2i}(S)$. 
 In particular, the following power series is well defined:
$$
\tilde \Theta(t):=\frac t {\Theta(t)}=1-a_2.t+(a_2^2-a_3).t^2+(-a_4+2a_2a_3-a^3_2).t^4+ \cdots
$$
The next result is a direct analog of the $\GL$-oriented case
 (see \cite[Prop. 4.1.2]{Deg12}).  
\end{num}

\begin{prop}
The Todd class $\td_\psi:\KSp_0 \rightarrow \F^{00 \times}$
 associated with the morphism $\psi$ (\Cref{prop:computeToddclass})
 is uniquely characterized by the relation 
$$
\td_\psi([\cU])=\left. \frac t {\Theta(t)}\right\vert_{t=b_1(\cU)}
$$
for a rank $2$ symplectic bundle $\cU$ over $X$.
\end{prop}
The case of an arbitrary symplectic bundle $\cU$ of rank $2n$ is determined
 by the symplectic splitting principle \Cref{num:sp-splitting-principle}.
 One introduces Borel roots
 $b_1^{(1)}=b_1(\cU_1)$, ..., $b_1^{(n)}=b_1(\cU_n)$ such that the Borel classes
 of $\cU$ satisfies the relations given by $b_t(\cU)=\prod_i (1+b_1^{(i)}.t)$.
 Then we obtain the expression
$$
\td_\psi([\cU])=\prod_{i=1}^n \tilde \Theta\big(b_1^{(i)}\big)
$$
which gives a computation of the left hand-side in terms of the Borel classes
 $b_i(\cU)$.

%\begin{rem}
% Moreover, \Cref{prop:sp-orient&FTL} implies that
% $(\Id,\phi):(\F^{**}(S),F^{b'}_t) \rightarrow (\F^{**}(S),F^{\phi(b')}_t)$ is a strict isomorphism of FTL.
%\end{rem}
